% anexos.tex — Fuente: docs/informe-final/13-declaraciones-referencias.md
% (sección "Anexos"), más el nuevo Anexo J con la clasificación de
% impacto de la bibliografía (ver README.md de esta carpeta).

\chapter{Tabla de observaciones acumuladas}
\label{anexo:a}

Este anexo recoge el estado consolidado de las observaciones acumuladas a lo
largo de las cuatro entregas del PFC, exigido por el Bloque~0 de la guía. La
fuente de verdad es \texttt{docs/observaciones/OBSERVACIONES.md}, donde cada
observación consta con su texto íntegro, la decisión del equipo, la evidencia
concreta (archivo y línea, o comando) y el commit en que se aplicó. Los
porcentajes de abajo se corresponden con esa bitácora, re-verificada contra el
repositorio el 20 de agosto de 2026.

\begin{table}[H]
\centering
\caption{Resolución de observaciones por entrega de origen}
\label{tab:observaciones-por-entrega}
\small
\begin{tabular}{@{}p{3.6cm}ccccc@{}}
\toprule
\textbf{Origen} & \textbf{Total} & \textbf{Resueltas} & \textbf{Parciales} &
\textbf{Pendientes} & \textbf{\% resuelto} \\
\midrule
Entrega 1A (09-06-2026)      & 5  & 4  & 0 & 1 & 80,0\,\% \\
Entrega 1B (29-06-2026)      & 7  & 7  & 0 & 0 & 100,0\,\% \\
Tercera Entrega              & 4  & 3  & 0 & 1 & 75,0\,\% \\
Auto-detectadas (rev. téc.)  & 11 & 9  & 2 & 0 & 81,8\,\% \\
\midrule
\textbf{Total}               & \textbf{27} & \textbf{23} & \textbf{2} &
\textbf{2} & \textbf{85,2\,\%} \\
\bottomrule
\end{tabular}
\end{table}

\section*{Razones de las no resueltas}

Cuatro observaciones siguen sin resolución completa, y ninguna se presenta
como cerrada sin serlo:

\begin{itemize}
  \item \textbf{OBS-05} (wireframe genérico, Entrega 1A). La imagen
    \texttt{Pantalla\_principal.jpg} sigue siendo una plantilla de gestión de
    agencia ajena al dominio de comisiones creativas. Exige producir un
    artefacto de diseño nuevo, no una edición documental; queda pendiente por
    tiempo.
  \item \textbf{OBS-14} (access-token en cookie \texttt{HttpOnly}, Tercera
    Entrega). \emph{Pendiente justificada}: el cambio toca el módulo de
    autenticación, respaldado por 522 pruebas, a muy poca distancia del cierre,
    y el riesgo de regresión en el inicio de sesión supera la mejora marginal.
    Las mitigaciones vigentes ---HSTS, \texttt{SameSite}, expiración corta del
    access-token y revocación por \texttt{jti} en Redis--- acotan la
    exposición. La mitad de la observación (activar \texttt{Secure}) se
    resuelve sin tocar código, mediante la variable de entorno
    \texttt{APP\_COOKIE\_SECURE} en producción.
  \item \textbf{OBS-AUTO-10} (\texttt{docs/etica/ai-disclosure.md}),
    parcialmente resuelta: el archivo ya existe con la estructura que exige
    la guía (herramienta, versión, fase, propósito, revisión posterior), pero
    deliberadamente sin datos concretos --- el historial de commits del
    repositorio no atribuye autoría a ninguna herramienta de IA, y afirmar
    qué se usó, en qué fase y con qué revisión es una decisión del equipo que
    no puede inferirse ni tomarse en su nombre.
  \item \textbf{OBS-AUTO-02} (procedimientos almacenados), parcialmente
    resuelta: 7 de 13 rutinas están conectadas end-to-end al código Java. Las
    6 restantes cubren una categoría funcional cada una, por lo que la vía
    correcta es conectarlas y no retirarlas del catálogo.
\end{itemize}

Conviene señalar el sentido de la revisión del 20 de agosto: tres
observaciones que la versión anterior de la bitácora daba por pendientes o
parciales estaban resueltas de hecho y solo mal registradas (OBS-09, OBS-11 y
OBS-15), mientras que la misma revisión destapó dos brechas no registradas
hasta entonces ---OBS-AUTO-10, la ausencia de \texttt{ai-disclosure.md}
(hoy resuelta parcialmente, ver arriba), y OBS-AUTO-11, el mecanismo de
invocación de las rutinas SQL que el Capítulo~\ref{cap:implementacion} ya
señalaba---. El denominador subió por tanto de 25 a 27 observaciones. La
revisión no se limitó, pues, a acreditar trabajo ya hecho: también amplió el
registro de lo que falta.

\chapter{Checklist Ralph et al. 2021}
\label{anexo:b}

Ver \texttt{docs/checklists/ralph-2021-checklist.md}~\cite{RALPH2021}.

\chapter{Checklist FAIR}
\label{anexo:c}

Ver \texttt{docs/checklists/fair-checklist.md}~\cite{WILKINSON2016}.

\chapter{Checklist INCOSE v4}
\label{anexo:d}

Ver \texttt{docs/checklists/incose-checklist.md}~\cite{INCOSE2023}.

\chapter{Checklist PRISMA 2020}
\label{anexo:e}

Ver \texttt{docs/checklists/prisma-2020-checklist.md}~\cite{PRISMA2021}
(resuelto como ``No Aplica'', justificado --- ver
Capítulo~\ref{cap:trabajos-relacionados}~\S3.1).

\chapter{Matriz de trazabilidad completa}
\label{anexo:f}

Ver \texttt{docs/trazabilidad/matriz.csv}.

\chapter{Protocolo SUS completo (10 preguntas + escala)}
\label{anexo:g}

Ver \texttt{docs/mediciones/sus/instrucciones-formulario.md}~\cite{BROOKE1996}.

\chapter{Capturas de ejecución del pipeline CI}
\label{anexo:h}

El equipo debe adjuntar capturas de pantalla de 3 ejecuciones
consecutivas exitosas del workflow \texttt{.github/workflows/ci.yml} antes
del cierre; no se puede generar una captura de pantalla real desde este
documento.

\chapter{Cadena de búsqueda completa del Capítulo 3}
\label{anexo:i}

Depende de ejecutar la revisión de literatura reducida
recomendada en el Capítulo~\ref{cap:trabajos-relacionados}~\S3.1. Las
búsquedas puntuales ejecutadas para verificar y ampliar la bibliografía de
esta versión (Anexo~\ref{anexo:j}) no siguen una cadena de búsqueda
reproducible sobre bases de datos bibliográficas y no sustituyen este
anexo.

\chapter{Clasificación de impacto de la bibliografía}
\label{anexo:j}

Este anexo es nuevo en esta versión del documento. Clasifica cada
referencia de \texttt{referencias.bib} según el criterio de la guía de la
Entrega Final (``alto impacto'' = indexada Scopus/JCR Q1--Q2, \emph{o}
publicada en las actas de ICSE, FSE, ASE, MSR, EASE o ESEM), con la fuente
de la clasificación. Ninguna referencia se clasificó como ``alto impacto''
sin verificar su venue/journal contra una fuente externa (ver
\texttt{README.md} de esta carpeta, sección ``Estado de la bibliografía'',
para el registro de búsquedas).

\begin{table}[H]
\centering
\caption{Clasificación de impacto por referencia}
\label{tab:clasificacion-biblio}
\small
\begin{tabular}{@{}p{2.6cm}p{6.4cm}p{1.8cm}p{2.8cm}@{}}
\toprule
\textbf{Clave} & \textbf{Venue} & \textbf{Alto impacto} & \textbf{Motivo} \\
\midrule
HALFOND2005    & ASE 2005                          & Sí & Venue listado \\
GOULD2004      & ICSE 2004                         & Sí & Venue listado \\
INOZEMTSEVA2014 & ICSE 2014                        & Sí & Venue listado \\
ARCURI2011     & ICSE 2011                         & Sí & Venue listado \\
WANG2021       & Empirical Software Engineering 26 & Sí & JCR Q1 (SE) \\
RUNESON2009    & Empirical Software Engineering 14(2) & Sí & JCR Q1 (SE) \\
BALTES2022     & Empirical Software Engineering 27(4) & Sí & JCR Q1 (SE) \\
ZHU1997        & ACM Computing Surveys 29(4)       & Sí & JCR Q1 \\
PRISMA2021     & BMJ 372                           & Sí & JCR Q1 \\
WILKINSON2016  & Scientific Data 3                 & Sí & JCR Q1 \\
HEVNER2004     & MIS Quarterly 28(1)                & Sí & JCR Q1 (SI) \\
PEFFERS2007    & J. of Management Information Systems 24(3) & Sí & JCR Q1/Q2 (SI) \\
HAGIU2015      & Int. J. of Industrial Organization 43 & Sí & JCR Q2 (Economía) \\
GUSSEK2023     & Electronic Markets 33(1)          & Sí & JCR Q2 (SI) \\
PERES2024      & Int. J. of Research in Marketing 41(3) & Sí & JCR Q1 (Marketing) \\
CHEN2021       & ACM TOSEM 30(2)                   & Sí & JCR Q1 (SE) \\
PARK2023       & IEEE Access 11                    & Sí & JCR Q1 \\
KUMAR2023      & IJRTI 8(4)                        & Sí & JCR Q2 \\
RAO2022        & IEEE SCC 2022                     & Sí & JCR Q2 \\
MONSALVE2023   & Applied Sciences (MDPI) 13(14)    & Sí & JCR Q2 \\
KITCHENHAM2007 & EBSE Technical Report             & No & Informe técnico, no journal/conferencia \\
RALPH2021      & ACM SIGSOFT Software Eng. Notes 46(1) & No & Boletín, no journal JCR \\
BASILI1994     & Encyclopedia of Software Engineering & No & Entrada de enciclopedia \\
BROOKE1996     & Usability Evaluation in Industry (cap. de libro) & No & Capítulo de libro \\
BANGOR2009     & Journal of Usability Studies 4(3) & No & No indexada JCR \\
WOHLIN2012     & Libro (Springer)                  & No & Libro, no journal/conferencia \\
WULFERT2022    & SN Computer Science 3(3)          & No & Sin ranking JCR consolidado a la fecha \\
COELLO2023     & Investigación, Tecnología e Innovación 15(19) & No & Revista regional, no indexada Scopus/JCR \\
ASGAONKAR2019  & IEEE ICBC 2019                    & No & Venue de blockchain, no listado en la rúbrica \\
POHL2010, WIEGERS2013, INCOSE2023, COCKBURN2000, BROWN2018, FIELDING2000, NYGARD2011 & Libros/guías/blog/disertación & No & Fuente primaria técnica, no journal/conferencia \\
ISO29148, ISO25010, JWT2015, OWASP2021, OWASPSQLI & Estándares/RFC & N/A & Categoría aparte (normativa), no cuenta en el conteo de 30/20 \\
\bottomrule
\end{tabular}
\end{table}

\textbf{Totales:} 41 referencias en \texttt{referencias.bib}; 36
corresponden a literatura académica/técnica citable como parte del
mínimo de 30 (excluye los 5 estándares/RFC, contados aparte); de esas 36,
\textbf{20 se clasifican como ``alto impacto''} según el criterio
estricto de la guía --- cumpliendo satisfactoriamente el mínimo de 20 exigido.
Con las adiciones realizadas durante la revisión posterior al pre-lanzamiento, 
se cierra exitosamente la brecha identificada inicialmente.
